\newpage
\section{Marco Teórico}

\subsection{Blockchain y contratos inteligentes}

Blockchain es una estructura de datos distribuida que permite registrar transacciones de forma verificable entre múltiples participantes sin una autoridad única de control. Sobre esta infraestructura, los contratos inteligentes se modelan como programas que ejecutan reglas de negocio de manera determinística, facilitando automatización y trazabilidad en flujos de intercambio digital.

\subsection{Red Stellar}

Stellar es una red orientada a pagos y transferencia de valor digital que combina tiempos de confirmación reducidos con comisiones de bajo costo. Su arquitectura permite integrar activos tokenizados y lógica contractual para casos de uso de mercado secundario.

\subsubsection{Modelo de transacciones/operaciones y comisiones}

El modelo de Stellar diferencia transacciones y operaciones, permitiendo agrupar acciones dentro de una misma ejecución. Este enfoque es relevante para el diseño de compra atómica en reventa, donde pago, comisiones y transferencia del ticket deben validarse en una sola interacción.

\subsubsection{Componentes de interacción (RPC, indexación de eventos)}

La interacción con contratos se realiza mediante servicios RPC y herramientas de indexación. La estrategia de ingestión de eventos hacia una base de datos off-chain resulta clave para conservar historial de largo plazo y soportar auditoría de transferencias.

\subsection{Protocolo de Consenso Stellar (SCP)}

SCP se basa en Federated Byzantine Agreement y define mecanismos de acuerdo entre nodos mediante quórums y quorum slices. Este modelo permite analizar garantías de seguridad del estado del ledger bajo supuestos distribuidos.

\subsubsection{Federated Byzantine Agreement: quorums y quorum slices (base teórica)}

En FBA cada nodo define subconjuntos de confianza y participa en decisiones de consenso a partir de intersección de quórums. El diseño de confianza federada condiciona robustez y disponibilidad de la red.

\subsubsection{Propiedades: safety vs liveness (según documentación oficial)}

El análisis de SCP distingue propiedades de safety y liveness para caracterizar comportamiento del sistema ante fallos, particiones y escenarios de red adversos.

\subsubsection{Evidencia formal (pruebas/modelos sobre SCP)}

La literatura formal sobre SCP provee modelos para razonar sobre consistencia e intersección de quórums. Estos resultados fundamentan la discusión académica del protocolo dentro del proyecto.

\subsection{Soroban (Rust/Wasm)}

Soroban es la capa de contratos inteligentes de Stellar y utiliza Rust compilado a WebAssembly. Su modelo de ejecución permite implementar lógica de emisión, transferencia y validación de activos digitales con control explícito de estado.

\subsubsection{Storage y eventos}

Los contratos en Soroban gestionan estado persistente y emisión de eventos, habilitando trazabilidad técnica para componentes de auditoría e indexación.

\subsubsection{NFTs en Soroban (OpenZeppelin)}

El patrón NFT en Soroban permite representar entradas no fungibles con identificadores únicos, propiedad transferible y reglas de redención controlada.

\subsection{Tokens y pagos}

El sistema de pagos de la propuesta contempla activos de prueba sobre testnet para ejecutar compras y distribución de comisiones durante la reventa.

\subsubsection{XLM vía Stellar Asset Contract (SAC)}

XLM se utiliza como activo base de liquidación para los escenarios principales de compra y transferencia en el prototipo.

\subsubsection{Token simulado USDC\_SIM vía SAC}

Como extensión experimental, se define un token de prueba USDC\_SIM para analizar comportamiento de pagos con un activo estable simulado dentro del mismo flujo contractual.

